\section{Security Analysis of API Commands}

\subsection{Proof Structure and Invariants}

We prove that each command preserves five flow invariants:

\begin{itemize}
    \item \textbf{I1}: stored objects have valid labels.
    \item \textbf{I2}: reads only to authorized principals.
    \item \textbf{I3}: writes only from authorized principals.
    \item \textbf{I4}: executed programs are flow-safe.
    \item \textbf{I5}: declassifications are authority-checked.
\end{itemize}

Assignment rule: $\Gamma,pc \vdash x := e$ iff $\Gamma \vdash e:L_e$ and $(pc \sqcup L_e \sqsubseteq \Gamma(x))$, preventing explicit and implicit flows.

\subsection{Command-by-Command Preservation}

\textbf{UPLOAD preserves I1:} Valid labels accepted; invalid ones rejected. No invalid label stored.

\textbf{READ\_DATA and WRITE\_DATA preserve I2/I3:} Both check access control before access. No unauthorized read/write.

\textbf{UPLOAD\_PROGRAM preserves I4:} Parser, type checker, and flow checker all reject unsafe code. No unsafe program stored.

\textbf{RUN\_PROGRAM preserves I4:} Executes only after flow-checking at upload time. Label check $L_{req} \sqsubseteq (outC, outI)$ ensures output label is high enough to hold all input information.

\textbf{DECLASSIFY\_AND\_RELEASE preserves I5:} Authority-checks that caller has required permissions for each relaxed policy, applies approved transform, and creates audit entry.

\textbf{AUDIT\_QUERY:} Read-only; preserves all invariants.


\subsection{Representative Illegal Flows and Blocking Points}

We present three realistic attack attempts and show where each is blocked:

\textbf{Case A: Explicit Read-Down Leak in User Code.} A researcher uploads a program:
\begin{verbatim}
  secret : (hospital:patients-only)
  public : (hospital:public-readable)
  ...
  public := secret
\end{verbatim}

The type checker verifies: $(hospital:\{patients\_only\}) \sqsubseteq (hospital:\{public\_readable\})$?

This requires $\{patients\_only\} \supseteq \{public\_readable\}$, which is false (the researcher is in $public\_readable$ but not in $patients\_only$). The flow check fails. UPLOAD\_PROGRAM rejects the code with a flow error. \textbf{Blocked at upload-time type checking.}

\textbf{Case B: Integrity Poisoning via Untrusted Input.} A researcher with low write authority attempts to run a program that modifies patient data with result to be stored with high-integrity label $(doctor \leftarrow \{doctor,hospital\})$:

RUN\_PROGRAM is called with:
\begin{itemize}
    \item input from researcher's low-integrity data: $L_{input} = (doctor \leftarrow \{researcher,hospital\})$,
    \item desired output label: $L_{out} = (doctor \leftarrow \{doctor,hospital\})$.
\end{itemize}

The check is: $L_{input} \sqsubseteq L_{out}$?

For integrity, this requires: for each $(o \leftarrow w_{out}) \in L_{out}$, there exists $(o \leftarrow w_{in}) \in L_{input}$ with $w_{in} \supseteq w_{out}$.

For $(doctor \leftarrow \{doctor,hospital\}) \in L_{out}$, we need $(doctor \leftarrow w) \in L_{input}$ with $w \supseteq \{doctor,hospital\}$. But $L_{input}$ has $(doctor \leftarrow \{researcher,hospital\})$ and $\{researcher,hospital\} \not\supseteq \{doctor,hospital\}$ (researcher is not trusted by the doctor). The check fails. \textbf{Blocked at run-time label check.}

\textbf{Case C: Unauthorized Downgrade via Declassification.} A hospital staff member (not a patient) attempts to declassify patient data from $(patient:\{patient\})$ to $(patient:\{researcher\})$:

DECLASSIFY\_AND\_RELEASE is called with:
\begin{itemize}
    \item source label: $(patient:\{patient\})$,
    \item target label: $(patient:\{researcher\})$,
    \item caller: $staff$ (not the patient).
\end{itemize}

The authority check: does the caller ($staff$) have authority from the patient to add researcher to the reader set?

This requires the caller to present a delegation or authorization token from the patient. If the caller is not the patient and has no such token, the authorization check fails. DECLASSIFY\_AND\_RELEASE rejects the call. \textbf{Blocked by authority check.}

\subsection{Trace-Level Theorem}

\textbf{Theorem (API Trace Security under Stated Assumptions).}
Assume:
\begin{enumerate}
    \item the secure authenticated transport layer (TLS with client and server authentication) is correctly implemented and prevents network-level adversaries from modifying commands or eavesdropping,
    \item the server's trusted enforcement code (the UPLOAD, READ, WRITE, UPLOAD\_PROGRAM, RUN\_PROGRAM, DECLASSIFY\_AND\_RELEASE, and AUDIT\_QUERY implementations) correctly execute the specified checks,
    \item user programs that have passed UPLOAD\_PROGRAM checking are executed in a runtime system that respects labeled data (i.e., does not perform implicit flows outside the labeled computation model).
\end{enumerate}

Then for any finite sequence (trace) of API calls that are all accepted by the server:
\begin{enumerate}
    \item all stored objects have valid labels satisfying I1,
    \item all data accessed via READ\_DATA respects confidentiality policies,
    \item all data modified via WRITE\_DATA respects integrity policies,
    \item all user programs executed via RUN\_PROGRAM have flow-safe interfaces matching their input and output labels,
    \item all declassifications performed via DECLASSIFY\_AND\_RELEASE are authority-checked and auditable,
\end{enumerate}

and the trace collectively preserves invariants I1--I5 with no illegal confidentiality or integrity flows except those explicitly permitted by authority-checked declassifications.

\textbf{Proof Sketch.} By induction on trace length: each command appended to a trace preserves the preceding invariants. Per-command arguments show each command type maintains I1--I5. Thus, any accepted trace preserves all invariants.
\begin{itemize}
    \item If $c_{n+1}$ is UPLOAD, I1 is preserved because label parsing and validation gate the command (shown above).
    \item If $c_{n+1}$ is READ\_DATA, I2 is preserved because the access control check gates the command.
    \item If $c_{n+1}$ is WRITE\_DATA, I3 is preserved because the integrity check gates the command.
    \item If $c_{n+1}$ is UPLOAD\_PROGRAM, I4 is preserved because type and flow checking gate the command.
    \item If $c_{n+1}$ is RUN\_PROGRAM, I4 is preserved because the label check gates the command.
    \item If $c_{n+1}$ is DECLASSIFY\_AND\_RELEASE, I5 is preserved because authority checking gates the command.
    \item If $c_{n+1}$ is AUDIT\_QUERY, no modification occurs, so invariants are unchanged.
\end{itemize}

By induction, any accepted trace satisfies all invariants. \qed

\textbf{Corollary: Non-Interference (Confidentiality).}
Consider two traces $\tau$ and $\tau'$ that differ only in the values of data objects labeled with owner $o$'s confidentiality policies. Let $P$ be a set of principals (readers) not including $o$. Then the sequence of responses and errors returned to $P$ along both traces is identical.

\emph{Intuition}: Principals not authorized by owner $o$ cannot observe differences in $o$'s protected data, because READ and DECLASSIFY both check authorization before revealing data.

\textbf{Limitations.}
The theorem does not cover:
\begin{itemize}
    \item compromised, malicious, or buggy server runtime or enforcement code,
    \item covert channels outside the labeled computation model (e.g., timing of responses, power consumption, acoustic noise),
    \item social engineering or the attacker forging authority credentials,
    \item misuse of legitimate authority by honest-but-curious insiders (the theorem prevents \emph{technical} leaks, but not intentional misuse of authorized operations).
\end{itemize}

These limitations are inherent to any formal security argument and do not diminish the value of the theorem in the specified threat model.

\subsection{Limitations}

The theorem does not cover:
\begin{itemize}
    \item compromised server runtime or keys,
    \item absence of all covert channels,
    \item legal/organizational misuse outside encoded policy checks.
\end{itemize}

These boundaries are explicitly reported so claims remain precise and defensible.
